%!TEX root = forallxyyc.tex
\chapter[快速参考]{快速参考}
%\pagestyle{plain}
\section{特征真值表}
\label{app.CharacteristicTTs}

\begin{tabular}{c|c}
\metav{A} & \enot\metav{A}\\
\hline
T & F\\
F & T \\
\phantom{.}\\
\phantom{.}
\end{tabular}


\hfill


\begin{tabular}{c c|c|c|c|c}
\metav{A} & \metav{B} & $\metav{A}\eand\metav{B}$ & $\metav{A}\eor\metav{B}$ & $\metav{A}\eif\metav{B}$ & $\metav{A}\eiff\metav{B}$\\
\hline
T & T & T & T & T & T\\
T & F & F & T & F & F\\
F & T & F & T & T & F\\
F & F & F & F & T & T
\end{tabular}

\vfill

\section{符号化}


\begin{center}
\label{app.symbolization}
\begin{tabular*}{\textwidth}{rl}
\multicolumn{2}{c}{\textsc{语句联结词}}\\ \\
并非 $P$ & $\enot P$\\
$P$ 或 $Q$ & $(P \eor Q)$\\
既非 $P$ 也非 $Q$ & $\enot(P \eor Q)$\ 或 \ $(\enot P \eand \enot Q)$\\
$P$ 且 $Q$ & $(P \eand Q)$\\
如果 $P$ 那么 $Q$ & $(P \eif Q)$\\
$P$ 仅当 $Q$ & $(P \eif Q)$\\
$P$ 当且仅当 $Q$ & $(P \eiff Q)$\\
$P$，除非 $Q$ & $(\enot Q \eif P)$ 或 $(P \eor Q)$\\
\\
\multicolumn{2}{c}{\label{SymbolizingPredicates}\textsc{谓词}}\\ \\
所有 $F$ 是 $G$ & $\forall x(\atom{F}{x} \eif \atom{G}{x})$\\
有些 $F$ 是 $G$ & $\exists x(\atom{F}{x} \eand \atom{G}{x})$\\
并非所有 $F$ 都是 $G$ & $\enot\forall x(\atom{F}{x} \eif \atom{G}{x})$\ 或\\
& $\exists x(\atom{F}{x} \eand \enot \atom{G}{x})$\\
没有 $F$ 是 $G$ & $\forall x(\atom{F}{x} \eif\enot \atom{G}{x})$\ 或\\
& $\enot\exists x(\atom{F}{x} \eand \atom{G}{x})$\\
只有 $F$ 才是 $G$ & $\forall x(\atom{G}{x} \eif \atom{F}{x})$\\
& $\enot\exists x(\enot\atom{F}{x} \eand \atom{G}{x})$\\
\\
\multicolumn{2}{c}{\textsc{等同}}\\ \\
只有 $c$ 是 $G$ & $\forall x(\atom{G}{x} \eiff x=c)$\\
除 $c$ 之外的一切\\都是 $G$ & $\forall x(\enot\, {x = c} \eif \atom{G}{x} )$\\
除 $c$ 之外的一切\\都是 $G$ & $\forall x(\enot\, {x = c} \eiff \atom{G}{x} )$\\
%$j$ is more $R$ than anyone else. & $\forall x(x\neq j \eif Rjx)$\\
那个 $F$ 是 $G$ & $\exists x(\atom{F}{x} \eand \forall y(\atom{F}{y} \eif x=y) \eand \atom{G}{x} )$\\
并非那个 $F$\\是 $G$ & $\enot\exists x(\atom{F}{x} \eand \forall y(\atom{F}{y} \eif x=y) \eand \atom{G}{x} )$\\
那个 $F$ 是非 $G$ & $\exists x(\atom{F}{x} \eand \forall y(\atom{F}{y} \eif x=y) \eand \enot \atom{G}{x} )$
\end{tabular*}
\end{center}

% BEGIN: symbolizing cardinality

\newpage
\section{使用等同来符号化数量}

\subsection*{至少有 \blank\ 个 $F$。}
\label{summary.atleast}

\ifHTMLtarget


\begin{tabular*}{\textwidth}{rl}
	一 & $\exists x\,\atom{F}{x}$\\
	二 & $\exists x_1\exists x_2(\atom{F}{x_1} \eand \atom{F}{x_2}
	\eand \enot x_1  = x_2)$\\
	三 & $\exists x_1\exists x_2\exists x_3(\atom{F}{x_1} \eand \atom{F}{x_2} \eand \atom{F}{x_3} \eand \enot x_1 = x_2 \eand\enot x_1 = x_3 \eand \enot x_2 = x_3)$\\
	四 & $\exists x_1\exists x_2\exists x_3\exists x_4 (\atom{F}{x_1} \eand \atom{F}{x_2} \eand \atom{F}{x_3} \eand \atom{F}{x_4} \eand \enot x_1 = x_2 \eand \enot x_1 = x_3 \eand \enot x_1 = x_4 \eand \enot x_2 = x_3 \eand \enot x_2 = x_4 \eand \enot x_3 = x_4)$\\
	$n$ & $\exists x_1\ldots\exists x_n(\atom{F}{x_1} \eand \ldots \eand \atom{F}{x_n} \eand \enot x_1 = x_2 \eand\ldots\eand \enot x_{n-1} = x_n)$
\end{tabular*}


\else


\begin{tabular*}{\textwidth}{rl}
	一 & $\exists x\,\atom{F}{x}$\\
	二 & $\exists x_1\exists x_2(\atom{F}{x_1} \eand \atom{F}{x_2}
	\eand \enot x_1  = x_2)$\\
	三 & $\exists x_1\exists x_2\exists x_3(\atom{F}{x_1} \eand \atom{F}{x_2} \eand \atom{F}{x_3} \eand {}$\\
	& \quad$\enot x_1 = x_2 \eand\enot x_1 = x_3 \eand \enot x_2 = x_3)$\\
	四 & $\exists x_1\exists x_2\exists x_3\exists x_4 (\atom{F}{x_1} \eand \atom{F}{x_2} \eand \atom{F}{x_3} \eand \atom{F}{x_4} \eand {}$\\
	& \quad $\enot x_1 = x_2 \eand \enot x_1 = x_3 \eand \enot x_1 = x_4 \eand {}$\\
	& \qquad$ \enot x_2 = x_3 \eand \enot x_2 = x_4 \eand \enot x_3 = x_4)$\\
	$n$ & $\exists x_1\ldots\exists x_n(\atom{F}{x_1} \eand \ldots \eand \atom{F}{x_n} \eand {}$\\
	& \quad$\enot x_1 = x_2 \eand\ldots\eand \enot x_{n-1} = x_n)$
\end{tabular*}


\fi

\subsection*{至多有 \blank\ 个 $F$。}
\label{summary.atmost}

表达“至多有 $n$ 个 $F$”的方式之一，是在“至少有 $n+1$ 个 $F$”的符号化前加上否定号。等价地，我们可以给出：

\noindent
\ifHTMLtarget


\begin{tabular*}{\textwidth}{rl}
	一 & $\forall x_1\forall x_2\bigl[(\atom{F}{x_1} \eand
	\atom{F}{x_2}) \eif x_1=x_2\bigr]$\\
	二 & $\forall x_1\forall x_2\forall x_3\bigl[(\atom{F}{x_1} \eand \atom{F}{x_2} \eand \atom{F}{x_3}) \eif (x_1=x_2 \eor x_1=x_3 \eor x_2=x_3)\bigr]$\\
	三 & $\forall x_1\forall x_2\forall x_3\forall x_4\bigl[(\atom{F}{x_1} \eand \atom{F}{x_2} \eand \atom{F}{x_3} \eand \atom{F}{x_4}) \eif (x_1=x_2 \eor x_1=x_3 \eor x_1=x_4 \eor x_2=x_3 \eor x_2=x_4 \eor x_3=x_4)\bigr]$\\
	$n$ & $\forall x_1\ldots\forall x_{n+1} \bigl[(\atom{F}{x_1} \eand \ldots \eand \atom{F}{x_{n+1}}) \eif (x_1=x_2 \eor \ldots \eor x_n=x_{n+1})\bigr]$
\end{tabular*}


\else


\begin{tabular*}{\textwidth}{rl}
	一 & $\forall x_1\forall x_2\bigl[(\atom{F}{x_1} \eand
	\atom{F}{x_2}) \eif x_1=x_2\bigr]$\\
	二 & $\forall x_1\forall x_2\forall x_3\bigl[(\atom{F}{x_1} \eand \atom{F}{x_2} \eand \atom{F}{x_3}) \eif {}$\\ & \quad$(x_1=x_2 \eor x_1=x_3 \eor x_2=x_3)\bigr]$\\
	三 & $\forall x_1\forall x_2\forall x_3\forall x_4\bigl[(\atom{F}{x_1} \eand \atom{F}{x_2} \eand \atom{F}{x_3} \eand \atom{F}{x_4}) \eif {}$\\
	& \quad$(x_1=x_2 \eor x_1=x_3 \eor x_1=x_4 \eor {}$\\
	& \qquad$x_2=x_3 \eor x_2=x_4 \eor x_3=x_4)\bigr]$\\
	$n$ & $\forall x_1\ldots\forall x_{n+1}
	\bigl[(\atom{F}{x_1} \eand \ldots \eand \atom{F}{x_{n+1}}) \eif {}$\\
	& \quad$(x_1=x_2 \eor \ldots \eor x_n=x_{n+1})\bigr]$
\end{tabular*}


\fi

\subsection*{恰好有 \blank\ 个 $F$。}
\label{summary.exactly}

表达“恰好有 $n$ 个 $F$”的一种方法是，将上述两种符号化用合取联结起来，说“至少有 $n$ 个 $F$ 且至多有 $n$ 个 $F$”。下列等价公式更为简短：

\noindent
\ifHTMLtarget


\begin{tabular*}{\textwidth}{rl}
	零 & $\forall x\,\enot \atom{F}{x}$\\
	一 & $\exists x\bigl[\atom{F}{x} \eand \forall y(\atom{F}{y} \eif x = y)\bigr]$\\
	二 & $\exists x_1\exists x_2\bigl[\atom{F}{x_1} \eand \atom{F}{x_2} \eand \enot x_1 = x_2 \eand \forall y\bigl(\atom{F}{y} \eif (y= x_1 \eor y = x_2)\bigr) \bigr]$\\
	三 & $\exists x_1\exists x_2\exists x_3\bigl[\atom{F}{x_1} \eand \atom{F}{x_2} \eand \atom{F}{x_3} \eand \enot x_1 =  x_2 \eand \enot  x_1 = x_3 \eand \enot x_2 = x_3 \eand \forall y\bigl(\atom{F}{y} \eif (y = x_1 \eor y = x_2 \eor y =  x_3)\bigr) \bigr]$\\
	$n$ & $\exists x_1\ldots\exists x_n\bigl[\atom{F}{x_1} \eand\ldots\eand \atom{F}{x_n}  \eand \enot x_1 = x_2 \eand\ldots\eand \enot x_{n-1}= x_n \eand \forall y\bigl(\atom{F}{y} \eif (y= x_1 \eor \ldots \eor y= x_n)\bigr)\bigr]$
\end{tabular*}


\else


\begin{tabular*}{\textwidth}{rl}
	零 & $\forall x\,\enot \atom{F}{x}$\\
	一 & $\exists x\bigl[\atom{F}{x} \eand \forall y(\atom{F}{y} \eif x = y)\bigr]$\\
	二 & $\exists x_1\exists x_2\bigl[\atom{F}{x_1} \eand \atom{F}{x_2} \eand {}$\\
	& \quad$\enot x_1 = x_2 \eand \forall y\bigl(\atom{F}{y} \eif (y= x_1 \eor y = x_2)\bigr) \bigr]$\\
	三 & $\exists x_1\exists x_2\exists x_3\bigl[\atom{F}{x_1} \eand \atom{F}{x_2} \eand \atom{F}{x_3} \eand {}$\\
	& \quad$\enot x_1 =  x_2 \eand \enot  x_1 = x_3 \eand \enot x_2 = x_3 \eand {}$\\
	& \qquad$\forall y\bigl(\atom{F}{y} \eif (y = x_1 \eor y = x_2 \eor y =  x_3)\bigr) \bigr]$\\
	$n$ & $\exists x_1\ldots\exists x_n\bigl[\atom{F}{x_1} \eand\ldots\eand \atom{F}{x_n}  \eand {}$\\
	& \quad$\enot x_1 = x_2 \eand\ldots\eand \enot x_{n-1}= x_n \eand {}$\\
	& \qquad$\forall y\bigl(\atom{F}{y} \eif (y= x_1 \eor \ldots \eor y= x_n)\bigr)\bigr]$
\end{tabular*}


\fi
%\item[one] $\exists x\forall y\bigl[\atom{F}{x} \eand (\atom{F}{y} \eif y = x)\bigr]$
%\item[two] $\exists x\exists y\forall z\Bigl(\atom{F}{x} \eand \atom{F}{y} \eand \bigl[\atom{F}{z} \eif (z=x \eor z=y)\bigr] \eand x \neq y\Bigr)$
%\item[three] $\exists x_1\exists x_2\exists x_3\forall y\Bigl(\atom{F}{x_1} \eand \atom{F}{x_2} \eand \atom{F}{x_3} \eand [\atom{F}{y} \eif (y=x_1 \eor y=x_2 \eor y=x_3)] \eand x_1 \neq x_2 \eand x_1 \neq x_3 \eand x_2 \neq x_3\Bigr)$
%\item[n] $\exists x_1\cdots\exists x_n\forall y\Bigl(\atom{F}{x_1} \eand \cdots \eand \atom{F}{x_n} \eand \bigl[\atom{F}{y} \eif (y=x_1 \eor \cdots \eor y=x_n)\bigr] \eand x_1 \neq x_2 \eand\cdots\eand x_{n-1}\neq x_n\Bigr)$

\label{ProofRules}
\newpage\section{真值函项逻辑基本推导规则}

% \mathem is an empty box the same width as $m$; used to align rules
% properly
\newlength{\mathemwd}
\settowidth{\mathemwd}{\mbox{\small$m$}}
\def\mathem{{\hbox to\mathemwd{}}}

% If compiling to PDF, make all the rules small(er) and avoid left margin
\ifHTMLtarget\else
\renewenvironment{fitchproof}
	{\noindent\par\noindent\small$\begin{nd}}
	{\end{nd}$\noindent\normalsize\ignorespacesafterend}
\fi

%{\LARGE \textbf{Basic Rules of Proof}}

\begin{multicols}{2}
\subsection*{重述}

\begin{fitchproof}
	\have[m]{a}{\metav{A}}
	\have[\ ]{c}{\metav{A}} \by{R}{a}
\end{fitchproof}

\subsection*{合取}

\begin{fitchproof}
	\have[m]{a}{\metav{A}}
	\have[n]{b}{\metav{B}}
	\have[\ ]{c}{\metav{A}\eand\metav{B}} \ai{a, b}
\end{fitchproof}
\begin{fitchproof}
	\have[m]{ab}{\metav{A}\eand\metav{B}}
	\have[\ ]{a}{\metav{A}} \ae{ab}
\end{fitchproof}
\begin{fitchproof}
	\have[m]{ab}{\metav{A}\eand\metav{B}}
	\have[\ ]{b}{\metav{B}} \ae{ab}
\end{fitchproof}

\subsection*{条件}

\begin{fitchproof}
	\open
		\hypo[m]{a}{\metav{A}}\AS
		\have[n]{b}{\metav{B}}
	\close
	\have[\ ]{ab}{\metav{A}\eif\metav{B}}\ci{a-b}
\end{fitchproof}
\begin{fitchproof}
	\have[m]{ab}{\metav{A}\eif\metav{B}}
	\have[n]{a}{\metav{A}}
	\have[\ ]{b}{\metav{B}} \ce{ab,a}
\end{fitchproof}

\subsection*{否定}

\begin{fitchproof}
\open
	\hypo[m]{a}{\metav{A}}\AS
	\have[n]{nb}{\ered}
\close
\have[\ ]{na}{\enot\metav{A}}\ni{a-nb}
\end{fitchproof}
\begin{fitchproof}
	\have[m]{na}{\enot\metav{A}}
	\have[n]{a}{\metav{A}}
	\have[ ]{bot}{\ered}\ri{na, a}
\end{fitchproof}

\subsection*{间接证明}

\begin{fitchproof}
\open
	\hypo[i]{a}{\enot\metav{A}}\AS
	\have[j]{nb}{\ered}
\close
\have[\mathem]{na}{\metav{A}}\ip{a-nb}
\end{fitchproof}


\subsection*{爆炸}

\begin{fitchproof}
\have[m]{bot}{\ered}
\have[ ]{}{\metav{A}}\re{bot}
\end{fitchproof}

\subsection*{析取}

\begin{fitchproof}
	\have[m]{a}{\metav{A}}
	\have[\ ]{ab}{\metav{A}\eor\metav{B}}\oi{a}
\end{fitchproof}
\begin{fitchproof}
	\have[m]{a}{\metav{A}}
	\have[\ ]{ba}{\metav{B}\eor\metav{A}}\oi{a}
\end{fitchproof}
\begin{fitchproof}
	\have[m]{ab}{\metav{A}\eor\metav{B}}
	\open
		\hypo[i]{a}{\metav{A}}\AS
		\have[j]{c1}{\metav{C}}
	\close
	\open
		\hypo[k]{b}{\metav{B}}\AS
		\have[l]{c2}{\metav{C}}
	\close
	\have[\mathem]{c}{\metav{C}} \oe{ab,a-c1, b-c2}
\end{fitchproof}

\subsection*{双条件}

\begin{fitchproof}
	\open
		\hypo[i]{a1}{\metav{A}} \AS
		\have[j]{b1}{\metav{B}}
	\close
	\open
		\hypo[k]{b2}{\metav{B}}\AS
		\have[l]{a2}{\metav{A}}
	\close
	\have[\mathem]{ab}{\metav{A}\eiff\metav{B}}\bi{a1-b1,b2-a2}
\end{fitchproof}
\begin{fitchproof}
	\have[m]{ab}{\metav{A}\eiff\metav{B}}
	\have[n]{a}{\metav{A}}
	\have[\ ]{b}{\metav{B}} \be{ab,a}
\end{fitchproof}
\begin{fitchproof}
	\have[m]{ab}{\metav{A}\eiff\metav{B}}
	\have[n]{a}{\metav{B}}
	\have[\ ]{b}{\metav{A}} \be{ab,a}
\end{fitchproof}

\end{multicols}

\section{真值函项逻辑的派生规则}


\begin{multicols}{2}
\subsection*{析取三段论}
\begin{fitchproof}
	\have[m]{ab}{\metav{A} \eor \metav{B}}
	\have[n]{nb}{\enot \metav{A}}
	\have[\ ]{con}{\metav{B}}\by{DS}{ab, nb}
\end{fitchproof}
\begin{fitchproof}
	\have[m]{ab}{\metav{A} \eor \metav{B}}
	\have[n]{nb}{\enot \metav{B}}
	\have[\ ]{con}{\metav{A}}\by{DS}{ab, nb}
\end{fitchproof}

\subsection*{否定后件}

\begin{fitchproof}
	\have[m]{ab}{\metav{A}\eif\metav{B}}
	\have[n]{a}{\enot\metav{B}}
	\have[\ ]{b}{\enot\metav{A}} \by{MT}{ab,a}
\end{fitchproof}

\subsection*{双重否定消去}
	\begin{fitchproof}
		\have[m]{dna}{\enot \enot \metav{A}}
		\have[ ]{a}{\metav{A}}\dne{dna}
	\end{fitchproof}


\subsection*{排中律}
	\begin{fitchproof}
		\open
			\hypo[i]{a}{\metav{A}}\AS
			\have[j]{c1}{\metav{B}}
		\close
		\open
			\hypo[k]{b}{\enot\metav{A}}\AS
			\have[l]{c2}{\metav{B}}
		\close
		\have[\mathem]{ab}{\metav{B}}\tnd{a-c1,b-c2}
	\end{fitchproof}

%
%\subsection*{Hypothetical Syllogism}
%
%\begin{fitchproof}
%	\have[m]{ab}{\metav{A}\eif\metav{B}}
%	\have[n]{bc}{\metav{B}\eif\metav{C}}
%	\have[\ ]{ac}{\metav{A}\eif\metav{C}}\by{HS}{ab,bc}
%\end{fitchproof}

\subsection*{德摩根规则}
\begin{fitchproof}
	\have[m]{ab}{\enot (\metav{A} \eor \metav{B})}
	\have[\ ]{dm}{\enot \metav{A} \eand \enot \metav{B}}\dem{ab}
\end{fitchproof}
\begin{fitchproof}
	\have[m]{ab}{\enot \metav{A} \eand \enot \metav{B}}
	\have[\ ]{dm}{\enot (\metav{A} \eor \metav{B})}\dem{ab}
\end{fitchproof}
\begin{fitchproof}
	\have[m]{ab}{\enot (\metav{A} \eand \metav{B})}
	\have[\ ]{dm}{\enot \metav{A} \eor \enot \metav{B}}\dem{ab}
\end{fitchproof}
\begin{fitchproof}
	\have[m]{ab}{\enot \metav{A} \eor \enot \metav{B}}
	\have[\ ]{dm}{\enot (\metav{A} \eand \metav{B})}\dem{ab}
\end{fitchproof}
\end{multicols}

%\newpage

\section{一阶逻辑的基本推导规则}

\begin{multicols}{2}
\subsection*{全称消去}

\begin{fitchproof}
	\have[m]{a}{\forall \metav{x}\,\metav{A}(\ldots \metav{x} \ldots \metav{x}\ldots)}
	\have[\ ]{c}{\metav{A}(\ldots \metav{c} \ldots \metav{c}\ldots)} \Ae{a}
\end{fitchproof}

\subsection*{全称引入}

\begin{fitchproof}
	\have[m]{a}{\metav{A}(\ldots \metav{c} \ldots \metav{c}\ldots)}
	\have[\ ]{c}{\forall \metav{x}\,\metav{A}(\ldots \metav{x} \ldots \metav{x}\ldots)} \Ai{a}
\end{fitchproof}
\begin{raggedright}
\metav{c} 不得出现在任何未撤销的假设中
\end{raggedright}
\end{multicols}

\begin{multicols}{2}
\subsection*{存在引入}

\begin{fitchproof}
	\have[m]{a}{\metav{A}(\ldots \metav{c} \ldots \metav{c}\ldots)}
	\have[\ ]{c}{\exists \metav{x}\,\metav{A}(\ldots \metav{x} \ldots \metav{c}\ldots)}\Ei{a}
\end{fitchproof}

\vfill\null\columnbreak

\subsection*{存在消去}

\begin{fitchproof}
	\have[m]{a}{\exists \metav{x}\,\metav{A}(\ldots \metav{x} \ldots \metav{x}\ldots)}
	\open
		\hypo[i]{b}{\metav{A}(\ldots \metav{c} \ldots \metav{c}\ldots)}\AS
		\have[j]{c}{\metav{B}}
	\close
	\have[\ ]{d}{\metav{B}}\Ee{a,b-c}
\end{fitchproof}
\begin{raggedright}
\noindent \metav{c} 不得出现在任何未解除的假设、$\exists \metav{x}\,\metav{A}(\ldots \metav{x} \ldots \metav{x}\ldots)$ 或 \metav{B} 中。
\end{raggedright}
\end{multicols}

\subsection*{同一性引入}

\begin{fitchproof}
	\have[\mathem]{x}{\metav{c}=\metav{c}} \by{=I}{}
\end{fitchproof}

\subsection*{同一性消去}

\begin{multicols}{2}
\begin{fitchproof}
	\have[m]{e}{\metav{a}=\metav{b}}
	\have[n]{a}{\metav{A}(\ldots \metav{a} \ldots \metav{a}\ldots)}
	\have[\ ]{ea1}{\metav{A}(\ldots \metav{b} \ldots \metav{a}\ldots)} \by{=E}{e,a}
\end{fitchproof}
\begin{fitchproof}
	\have[m]{e}{\metav{a}=\metav{b}}
	\have[n]{a}{\metav{A}(\ldots \metav{b} \ldots \metav{b}\ldots)}
	\have[\ ]{ea2}{\metav{A}(\ldots \metav{a} \ldots \metav{b}\ldots)} \by{=E}{e,a}
\end{fitchproof}
\end{multicols}

%\begin{minipage}{\textwidth} % hack to keep section header with table
\section{一阶逻辑的派生规则}

\begin{multicols}{2}
\begin{fitchproof}
	\have[m]{ab}{\forall \metav{x}\enot \metav{A}}
	\have[\ ]{ac}{\enot \exists \metav{x} \metav{A}}\cq{m}
\end{fitchproof}
\begin{fitchproof}
	\have[m]{ab}{\enot \exists \metav{x}  \metav{A}}
	\have[\ ]{ac}{\forall \metav{x}\enot\metav{A}}\cq{m}
\end{fitchproof}
\begin{fitchproof}
	\have[m]{ab}{\exists \metav{x}\enot\metav{A}}
	\have[\ ]{ac}{\enot \forall \metav{x} \metav{A}}\cq{m}
\end{fitchproof}
\begin{fitchproof}
	\have[m]{ab}{\enot \forall \metav{x}  \metav{A}}
	\have[\ ]{ac}{\exists \metav{x}\enot \metav{A}}\cq{m}
\end{fitchproof}
\end{multicols}


%\end{minipage}

\section{子证明与引用的使用规则}

在应用规则时，若要引用单独一行：

\begin{compactlist}
\item 该行必须出现在规则应用行之前，并且
\item 不得出现在一个于规则应用行之前就已闭合的子证明之内。
\end{compactlist}

在应用规则时，若要引用一个子证明：

\begin{compactlist}
\item 被引用的子证明必须完全出现在引用它的规则应用之前，
\item 被引用的子证明不得位于某个在引用行处已经闭合的其他已闭合子证明之内，并且
\item 被引用子证明的最后一行不得位于一个嵌套子证明的内部。
\end{compactlist}

\section{等价链规则}

\ifHTMLtarget
\begin{align*}
\lnot\lnot \metav{P} & \Leftrightarrow \metav{P} && \text{(DN)}\\[2ex]
(\metav{P} \eif \metav{Q}) & \Leftrightarrow (\lnot \metav{P} \lor \metav{Q})
&& \text{(Cond)}\\
\lnot(\metav{P} \eif \metav{Q}) & \Leftrightarrow (\metav{P} \land \lnot\metav{Q}) \\[2ex]
(\metav{P} \eiff \metav{Q}) & \Leftrightarrow ((\metav{P} \eif \metav{Q}) \land  (\metav{Q} \eif \metav{P}))
&& \text{(Bicond)}\\[2ex]
\lnot(\metav{P} \land \metav{Q}) & \Leftrightarrow (\lnot\metav{P} \lor \lnot\metav{Q})
&& \text{(DeM)}\\
\lnot(\metav{P} \lor \metav{Q}) & \Leftrightarrow (\lnot\metav{P} \land \lnot\metav{Q}) \\[2ex]
(\metav{P} \lor \metav{Q}) & \Leftrightarrow (\metav{Q} \lor \metav{P}) & &\text{(Comm)}\\
(\metav{P} \land \metav{Q}) & \Leftrightarrow (\metav{Q} \land \metav{P})\\[2ex]
(\metav{P} \land (\metav{Q} \lor \metav{R})) & \Leftrightarrow ((\metav{P} \land \metav{Q}) \lor (\metav{P} \land \metav{R}))
&& \text{(Dist)}\\
(\metav{P} \lor (\metav{Q} \land \metav{R})) & \Leftrightarrow ((\metav{P} \lor \metav{Q}) \land (\metav{P} \lor \metav{R}))\\
(\metav{P} \lor (\metav{Q} \lor \metav{R})) & \Leftrightarrow ((\metav{P} \lor \metav{Q}) \lor \metav{R}) & &\text{(Assoc)}\\
(\metav{P} \land (\metav{Q} \land \metav{R})) & \Leftrightarrow ((\metav{P} \land \metav{Q}) \land \metav{R})\\[2ex]
(\metav{P} \lor \metav{P}) & \Leftrightarrow \metav{P} && \text{(Id)}\\
(\metav{P} \land \metav{P}) & \Leftrightarrow \metav{P}\\[2ex]
(\metav{P} \land (\metav{P} \lor \metav{Q})) & \Leftrightarrow \metav{P}&& \text{(Abs)}\\
(\metav{P} \lor (\metav{P} \land \metav{Q})) & \Leftrightarrow \metav{P}\\[2ex]
(\metav{P} \land (\metav{Q} \lor \lnot\metav{Q})) & \Leftrightarrow \metav{P}  && \text{(Simp)}\\
(\metav{P} \lor (\metav{Q} \land \lnot\metav{Q})) & \Leftrightarrow \metav{P}\\[2ex]
(\metav{P} \lor (\metav{Q} \lor \lnot\metav{Q})) & \Leftrightarrow (\metav{Q} \lor \lnot\metav{Q}) \\
(\metav{P} \land (\metav{Q} \land \lnot\metav{Q})) & \Leftrightarrow (\metav{Q} \land \lnot\metav{Q})
\end{align*}
\else

\begin{align*}
\lnot\lnot \metav{P} & \Leftrightarrow \metav{P} && \text{(DN)}\\[2ex]
(\metav{P} \eif \metav{Q}) & \Leftrightarrow (\lnot \metav{P} \lor \metav{Q})
&& \text{(Cond)}\\
\lnot(\metav{P} \eif \metav{Q}) & \Leftrightarrow (\metav{P} \land \lnot\metav{Q}) \\[2ex]
(\metav{P} \eiff \metav{Q}) & \Leftrightarrow ((\metav{P} \eif \metav{Q}) \land  (\metav{Q} \eif \metav{P}))
&& \text{(Bicond)}\\[2ex]
\lnot(\metav{P} \land \metav{Q}) & \Leftrightarrow (\lnot\metav{P} \lor \lnot\metav{Q})
&& \text{(DeM)}\\
\lnot(\metav{P} \lor \metav{Q}) & \Leftrightarrow (\lnot\metav{P} \land \lnot\metav{Q}) \\[2ex]
(\metav{P} \lor \metav{Q}) & \Leftrightarrow (\metav{Q} \lor \metav{P}) & &\text{(Comm)}\\
(\metav{P} \land \metav{Q}) & \Leftrightarrow (\metav{Q} \land \metav{P})\\[2ex]
(\metav{P} \land (\metav{Q} \lor \metav{R})) & \Leftrightarrow ((\metav{P} \land \metav{Q}) \lor (\metav{P} \land \metav{R}))
&& \text{(Dist)}\\
(\metav{P} \lor (\metav{Q} \land \metav{R})) & \Leftrightarrow ((\metav{P} \lor \metav{Q}) \land (\metav{P} \lor \metav{R}))\\
(\metav{P} \lor (\metav{Q} \lor \metav{R})) & \Leftrightarrow ((\metav{P} \lor \metav{Q}) \lor \metav{R}) & &\text{(Assoc)}\\
(\metav{P} \land (\metav{Q} \land \metav{R})) & \Leftrightarrow ((\metav{P} \land \metav{Q}) \land \metav{R})\\[2ex]
(\metav{P} \lor \metav{P}) & \Leftrightarrow \metav{P} && \text{(Id)}\\
(\metav{P} \land \metav{P}) & \Leftrightarrow \metav{P}\\[2ex]
(\metav{P} \land (\metav{P} \lor \metav{Q})) & \Leftrightarrow \metav{P}&& \text{(Abs)}\\
(\metav{P} \lor (\metav{P} \land \metav{Q})) & \Leftrightarrow \metav{P}\\[2ex]
(\metav{P} \land (\metav{Q} \lor \lnot\metav{Q})) & \Leftrightarrow \metav{P}  && \text{(Simp)}\\
(\metav{P} \lor (\metav{Q} \land \lnot\metav{Q})) & \Leftrightarrow \metav{P}
\end{align*}
\begin{align*}
(\metav{P} \lor (\metav{Q} \lor \lnot\metav{Q})) & \Leftrightarrow (\metav{Q} \lor \lnot\metav{Q}) \\
(\metav{P} \land (\metav{Q} \land \lnot\metav{Q})) & \Leftrightarrow (\metav{Q} \land \lnot\metav{Q})
\end{align*}
\fi
